\section*{NS90: a signed M\"obius attempt on the PR42 lower gain}

\paragraph{Scope.}
This is an exact arithmetic identity, a proof about the loss from one
absolute-value estimate, and a certified finite test of that estimate.
It is not a lower-gain theorem. The coefficient direction itself was
already tested in NS73; its finite efficiency is not a new result here.
The full joint observation energy of NS87 remains the principal open target.

Work in $L^2((0,\infty),dt/t^2)$, with
$a_n(t)=A(t/n)$, $A(z)=\int_0^z\{u\}\,du/u$, and
$\tau(t)=\mathbf1_{t\ge1}\log t$.
Write $R=\tau-\Pi_N\tau$, $E=\|R\|^2$, $h_n=\langle R,a_n\rangle$;
then $h_n=0$ for $n\le N$ and $\langle R,\tau\rangle=E$.

\paragraph{Exact identity: the prime main terms cancel.}
Define the unnormalized classical logarithmic taper and its residual by
\[
 F_X^\mu=-\sum_{n\le X}\mu(n)\log(X/n)a_n,
 \qquad Q_X=(\log X)\tau-F_X^\mu,
 \qquad S_X=\sum_{n\le X}\frac{\mu(n)}n\log(X/n).
\]
The divisor identity of NS85, now keeping every omitted divisor, gives
\[
 Q_X(t)=S_Xt-\Psi_1(t)+\sum_{X<k\le t}b_X(k)\log(t/k),
 \qquad b_X(k)=\sum_{\substack{d\mid k\\d>X}}\mu(d)\log(X/d),
 \qquad \Psi_1(t)=\sum_{k\le t}\Lambda(k)\log(t/k).
\]
For $t<1$ the two sums vanish. This follows from
$\sum_{d\mid k}\mu(d)=\mathbf1_{k=1}$ and
$\sum_{d\mid k}\mu(d)\log d=-\Lambda(k)$.
At integer endpoints the newly entering logarithmic factor is zero.

Set $V=F_{2N}^\mu-F_N^\mu$, $D=Q_{2N}-Q_N=(\log2)\tau-V$.
The old normal equations remove the coefficients of $V$ below $N$:
\begin{equation}\label{ns90-numerator}
 \ell_N:=\sum_{N<n\le2N}-\mu(n)\log(2N/n)h_n
       =(\log2)E-\langle R,D\rangle.
\end{equation}
The two identical $-\Psi_1$ terms cancel exactly in $D$.
Thus they cannot be counted as an independent favorable main term in
\eqref{ns90-numerator}.
Explicitly, put $s=S_{2N}-S_N$ and
\[
 d_n=\begin{cases}-\mu(n)\log2&n\le N,\\
                 -\mu(n)\log(2N/n)&N<n\le2N,\end{cases}
 \qquad \delta_k=(\log2)\mathbf1_{k=1}+\sum_{n\mid k}d_n.
\]
Then $s=-\sum d_n/n$, $\delta_k=0$ through $N$, and
$\delta_k=\mu(k)\log(k/N)$ for $N<k\le2N$.
For all larger $k$ the complete divisor sum must be retained. Globally,
\[
 D(t)=st+\sum_{k\le t}\delta_k\log(t/k).
\]

\paragraph{Exact finite formula and complete tail.}
For $T\ge2N$, define
\[
 I_T=\int_0^T R(t)\frac{dt}{t},\qquad
 P_T(k)=\int_k^T\log(t/k)R(t)\frac{dt}{t^2},\qquad
 Z_T=\int_T^\infty R(t)D(t)\frac{dt}{t^2}.
\]
Finite interchange, including $0<t<1$, yields
\begin{equation}\label{ns90-finite}
 \ell_N=(\log2)E-sI_T-\sum_{N<k\le T}\delta_kP_T(k)-Z_T.
\end{equation}
For a residual $r=\alpha\tau-\sum c_na_n$, the complete NS67 remainder is
\[
 r(t)=u\log t+v+\epsilon(t),\quad
 u=\alpha-\tfrac12\sum c_n,\quad
 v=\tfrac12\sum c_n\log(n/(2\pi)),\quad
 |\epsilon(t)|\le B/t,\quad B=\tfrac16\sum n|c_n|
\]
when $t$ exceeds its coefficient support. With $q_i=u_i\log T+v_i$,
\[
 Z_0=\frac{q_Rq_D+u_Rq_D+u_Dq_R+2u_Ru_D}{T},\quad
 M_i=\frac{q_i^2+2u_iq_i+2u_i^2}{T},\quad U_i=\frac{B_i^2}{3T^3},
\]
\[
 |Z_T-Z_0|\le\sqrt{M_RU_D}+\sqrt{M_DU_R}+\sqrt{U_RU_D}.
\]
The two mixed terms are essential. All these are complete-space formulas.

\paragraph{Proved implication: why termwise absolute values can get worse.}
Fix $N\ge2$. If the actual residual's pair $(u_R,v_R)$ is nonzero, then
\begin{equation}\label{ns90-absolute}
 \sum_{N<k\le T}|\delta_kP_T(k)|\longrightarrow\infty
 \quad (T\longrightarrow\infty).
\end{equation}
Here $N$ is fixed; this is not a uniform statement at $T=N^6$ as $N$ grows.
To prove it, let $L=\operatorname{lcm}(1,\ldots,2N)$.
For $k>1$, $\delta_k$ is periodic modulo $L$.
At $k=1+jL$, $j\ge1$, no $n\in\{2,\ldots,2N\}$ divides $k$;
hence $\delta_k=d_1=-\log2\ne0$ on this entire progression.
The complete potential has the asymptotic, with explicit error,
\[
 P_\infty(k)=\frac{u_R\log k+2u_R+v_R}{k}+O(B_R/(4k^2)).
\]
If $u_R\ne0$, its absolute sum on this progression diverges like
$(\log X)^2$; if $u_R=0$, $v_R\ne0$, it diverges like $\log X$.
For $k\le\sqrt T$, the uniform tail bound
$|P_\infty(k)-P_T(k)|=O_{R}(\log^2T/T)$ shows that the total replacement
error on that progression is $O_R(\log^2T/\sqrt T)=o(1)$.
Restrict the nonnegative sum to these indices to obtain
\eqref{ns90-absolute}. No prime-distribution asymptotic or RH input is used.

Consequently the elementary lower bound
\[
 L_T=(\log2)E-|sI_T|-\sum_{N<k\le T}|\delta_kP_T(k)|-|Z_T|
\]
tends to minus infinity under the stated fixed-$N$ hypothesis, even though
the signed formula converges. It is the estimate that loses the
cancellation, not an obstruction to the signed direction or full joint
route. The finite certificates verify $u_R\ne0$ at the two tested sizes.
The asymptotic argument does not provide uniform constants in $N$ or
exclude a different bound at a coupled cutoff.

\paragraph{Certified finite test.}
The original optimized coefficient balls from NS89, not its normalized
rational trial vectors, are integrated at $N=16,256$, $T=262144$, using
256 and 384 bits. The $N=256$ replay doubles both the inherited Gram
series cutoff and $T$. Every complete numerator overlaps the NS73
full-Gram numerator and is positive. At $N=256$ and $T=262144$,
$L_T/E<-186$; after doubling $T$, $L_T/E<-226$.
At $N=16$ the first bound is below $-828E$.
Dropping the finite divisor forcing gives a negative numerator at both
sizes, whereas the complete numerator is positive. This diagnoses that
particular omission; it is not an asymptotic lower estimate.

\paragraph{Open input and actual cost.}
With $W=(I-\Pi_N)V$ and $K_N=\|W\|^2$, exact line search gives
$E_N-E_{2N}\ge\ell_N^2/K_N$ when $K_N>0$.
The values of $K_N$ recorded here are the complete projected costs already
certified in NS73; no raw or model cost is substituted. A useful cofinal
argument still requires a one-sided signed estimate in
\eqref{ns90-finite} together with a compatible true-cost bound, or a direct
lower estimate for NS87's full joint observation energy. None is proved.
In particular the finite positive numerators do not establish the needed
nonsummable relative gain. RH, G2 and the original uniform Weil floor
remain open. No manuscript version is added for this attempt.
